\begin{summarybox}
    \textbf{\textcolor{CSummary}{一句话核心}}

    对$p>0$，若$P(x_0,y_0)$在抛物线$y^2=2px$上，即$y_0^2=2px_0$，
    则$P$处切线为
    \[
        \boxed{y_0y=p(x+x_0)}.
    \]
    \textbf{理解线索：}两点相减得到割线的方向，割线靠近切点得到切线；
    $x+x_0$来自$y_0^2=2px_0$的整理。

    \medskip
    \textbf{\textcolor{CSummary}{三种证明：按使用门槛排列}}
    \begin{itemize}[leftmargin=*, itemsep=0.4em, topsep=0.3em]
        \item \textbf{（1） 最低门槛（平方差）：}
        $(y_1-y_0)(y_1+y_0)=2p(x_1-x_0)$，割线靠近切线，
        得$x-x_0=\dfrac{y_0}{p}(y-y_0)$。
        \item \textbf{（2） 证明最快（导数）：}
        $x=f(y)=\dfrac{y^2}{2p}$，求导得$f'(y)=\dfrac{y}{p}$，
        直接写出切线。
        \item \textbf{（3） 通用证明（判别式）：}
        设$x=my+b$，代入曲线并令$\Delta=0$，得
        $y_0=pm$、$b=-x_0$，从而得到切点式。
    \end{itemize}

    \medskip
    \textbf{\textcolor{CSummary}{新增推论：切线轴上截距与中点}}

    设切点$P(x_0,y_0)$，$PH\perp x$轴于$H$，切线交$x$轴于$T$。
    把$y=0$代入切点式，可得
    \[
      \boxed{T(-x_0,0)},\qquad H(x_0,0),\qquad
      \boxed{O\text{是}TH\text{的中点}\ (P\ne O)}.
    \]
    \textbf{考场速用：}切线在$x$轴上的截距$t=-x_0$，反过来$x_0=-t$。
    当$P=O$时，$T=H=O$，仍有$t=-x_0=0$。
    对$x^2=\pm2py$，对应改为$y$轴截距$-y_0$。

    \medskip
    \textbf{\textcolor{CSummary}{使用条件与考场提醒}}
    \begin{itemize}[leftmargin=*, itemsep=0.4em, topsep=0.3em]
        \item 切点必须\textbf{在抛物线上}；切线经过的任意已知点不一定是切点。
        \item $y_0\ne0$时切线斜率$k=\dfrac{p}{y_0}$；$y_0=0$时顶点切线为$x=0$。
        \item 其他开口方向必须换成对应公式；\textbf{一个公共点不一定说明相切}。
    \end{itemize}
\end{summarybox}
